\section{座標環による実現と補完定理の射程}

量子多項式環
\[
 S=k\langle x_1,x_2,x_3,w\rangle/
 (x_jx_i-q_{ij}x_ix_j,\ wx_i-\lambda_i x_iw),
 \qquad q_{ij},\lambda_i\ne0
\]
は、重み $(1,1,1,3)$ のHilbert級数
$((1-z)^3(1-z^3))^{-1}$ を持つ。
この場合はAuroux--Katzarkov--Orlovの定理をそのまま適用でき、
$S_{(0,1,3,4)}$ のperfect圏は $D^b(\operatorname{qgr}S)$ の許容部分圏となる。
別稿\cite{VKPoisson32}では、二変数四次式 $f$ に対し
\[
 \frac{k\langle y,z,x,w\rangle}{
 (yz-zy,\ xy-ayx,\ xz-bzx,\ wy-a^{-1}yw,\ wz-b^{-1}zw,\ wx-abxw+f)}
\]
という反復Ore拡大を構成する。
$a,b\ne0$ のとき、これはAS正則次元四、Gorenstein parameter六の4-Calabi--Yau代数である。
Orlovの定理\cite[Corollary 2.7]{Orlov2009}により標準六項のfull strong列を持つため、
この族についても実際の $\Perf(S_{(0,1,3,4)})$ の半直交実現が得られる。

一方、本稿の一般量子化定理が直接構成するのは、各次数成分が自由である形式的な重み付き代数と、
その六項の有限切り出しである。
任意の形式量子化に対するAS正則性、パラメータの代数化、あるいは解析的収束性はここでは使わない。
従って主定理の「周囲」は、明確に指定した $\Lambda$ の形式変形として理解する。
全ての非可換重み付き射影空間の分類という意味には広げない。

既存の比較で現れる
\[
 \operatorname{gr}_J B_{p,q}\simeq(Q_{p,q})_{(0,1,3,4)}
\]
は量子中間代数、すなわち随伴次数代数の実現を与える。
この等式だけから元の $B_{p,q}$ 自身の埋め込みは従わない。
本稿の $st=0$ 上のOre族では、随伴次数を取ることなく実際のcornerを同定している。

また、次数 $1,1,1,3$ の最小生成元と上のHilbert級数を持つ座標環では、
\[
 \dim S_3/(S_1S_2)=1,\qquad
 \operatorname{rank}(S_1\otimes S_2\longrightarrow S_3)=10
\]
である。したがって、この合成が階数十一である標準四項代数を、そのまま
$(0,1,3,4)$ の標準部分列と同定することはできない。
この階数制約は、標準生成対象を変更した導来埋め込みの不存在を意味しない。


\subsection{残る問題の区別}
Poissonスキームの少なくとも八成分と、三角形・共点三直線・共鳴点の局所環は別稿で決定した。
そのスキームは非被約である。しかし全成分の分類と最大次元は残る。
一般のPoisson量子化をnoetherian AS正則座標環で代数化する問題も、有限六項の形式補完とは別である。

別の例外列まで許した埋め込み問題には、本稿の標準補完障害を越える不変量が必要となる。
一方、同じ標準cornerに対する補完の同型類が高々一つであることは、
新しい幾何的実現を得たとき、その実現を既存の六項族と比較する道具となる。

\section{計算の再現範囲}

同梱のPython検証器は標準ライブラリの整数と有理数だけで実行できる。
通常実行は証明書を変更せず、再計算結果を保存済み証明書と比較する。
統合した入口は \texttt{code/verify\_v31.py} である。
Hochschild複体の微分、制限のchain map、評価行列の階数、Schouten式の全係数、
BVホモトピー、二変数族の全ての非恒等三重合成を検査する。

Poissonスキームの新規成分と局所環の検証は別稿の六つのゲートにまとめ、
\nolinkurl{code/verify_v32.py} から実行できる。

全Artin基底上の補完条件、全次数の量子化、二次的miniversal基底の同定は、有限計算の射程には含めない。
それぞれcorner同型の持上げ、Kontsevichの定理、二次生成とNakayamaの補題が証明の根拠である。
